% !TeX program = tectonic
\documentclass[11pt,a4paper]{article}
\usepackage{fontspec}
\usepackage[russian]{babel}
\babelfont{rm}[Language=Default]{Liberation Serif}
\babelfont[Language=Default]{sf}{Carlito}
\babelfont[Language=Default]{tt}{DejaVu Sans Mono}
\usepackage{amsmath,amssymb,amsthm}
\usepackage{geometry}
\geometry{a4paper, top=2.4cm, bottom=2.4cm, left=2.4cm, right=2.4cm}
\usepackage{booktabs,tabularx,enumitem,xcolor,tcolorbox}
\tcbuselibrary{breakable,skins}
\definecolor{linkblue}{HTML}{1F4E79}
\definecolor{statgreen}{HTML}{2F6B3A}
\definecolor{goldacc}{HTML}{8B7E5A}
\usepackage[colorlinks=true,linkcolor=linkblue,urlcolor=linkblue,unicode=true]{hyperref}
\theoremstyle{plain}
\newtheorem{theorem}{Теорема}
\newtheorem{lemma}[theorem]{Лемма}
\newtheorem{proposition}[theorem]{Предложение}
\newtheorem{corollary}[theorem]{Следствие}
\theoremstyle{definition}
\newtheorem{definition}[theorem]{Определение}
\theoremstyle{remark}
\newtheorem{remark}[theorem]{Замечание}
\newtcolorbox{statusbox}[1][]{colback=statgreen!5,colframe=statgreen!75,
  title={\textbf{Сводка доказанного:} #1},fonttitle=\small,boxrule=0.7pt,arc=2pt,breakable}
\newtcolorbox{databox}[1][]{colback=goldacc!6,colframe=goldacc!80,
  title={\textbf{Данные протокола:} #1},fonttitle=\small,boxrule=0.7pt,arc=2pt,breakable}
\newtcolorbox{verifybox}[1][]{colback=linkblue!5,colframe=linkblue!80,
  title={\textbf{Верификация:} #1},fonttitle=\small,boxrule=0.7pt,arc=2pt,breakable}
\widowpenalty=10000
\clubpenalty=10000
\setlength{\parskip}{0.35em}
\newcommand{\CC}{\mathbb{C}}
\newcommand{\RR}{\mathbb{R}}
\newcommand{\ZZ}{\mathbb{Z}}
\newcommand{\QQ}{\mathbb{Q}}
\newcommand{\Ga}{\Gamma}
\newcommand{\Bfun}{\mathrm{B}}
\newcommand{\im}{\operatorname{Im}}
\newcommand{\re}{\operatorname{Re}}
\newcommand{\lcm}{\operatorname{lcm}}
\newcommand{\SNF}{\operatorname{SNF}}
\newcommand{\Jac}{\operatorname{Jac}}
\newcommand{\rank}{\operatorname{rank}}
\newcommand{\Ind}{\operatorname{Index}}
\newcommand{\disc}{\operatorname{disc}}
\begin{document}
\begin{center}
{\LARGE\bfseries Сертификат H и стенд Клейна}\\[0.4em]
{\large Γ-нормировка Гросса, объём 1125 и арифметика гептады}\\[0.6em]
{\normalsize Исаев Исхак Хамзатович}\\[0.2em]
{\small Программа <<Динамический принцип>> --- лаборатория \texttt{hodge-laboratory}}\\[0.15em]
{\small Отдельная монография теоремы · Монография 16/16}
\end{center}
\vspace{0.4em}
{\color{linkblue}\hrule height 0.8pt}
\vspace{0.8em}

\section{Постановка}

Сертификат H --- вершина арки программы: полная $\Ga$-нормировка
Гросса на стенде $N=15/30$, закрывающая $\Ga$-остаток сертификата G.
Вместе с ним здесь собран и третий стенд --- якобиан квартики
Клейна с его арифметикой.

\begin{theorem}[H1 --- циклотомическая решётка]\label{th:H1stand}
CM-решётка стенда: $\QQ(\zeta_{30})=\QQ(\zeta_{15})$ --- один
стенд, два $\Ga$-уровня. Форма Римана целочисленна (суммы
Рамануджана); тип поляризации (SNF) $(1,1,5,5,15,15,15,15)$ ---
целые индексы; $\disc=3^4\cdot5^6=1265625=1125^2$.
\end{theorem}

\begin{proof}
Форма Рамануджа целочисленна; SNF вычислена точным целочисленным
алгоритмом; произведение инвариантных множителей
$1\cdot1\cdot5\cdot5\cdot15^4=1265625$ совпадает с $3^4\cdot5^6$
--- разложение проверено. Поле: $\phi(30)=\phi(15)=8$;
изоморфизм $\zeta_{30}\mapsto-\zeta_{15}^8$.
\end{proof}

\begin{theorem}[H2 --- $\Ga$-нормировка объёма]\label{th:H2stand}
\[
  \mathrm{vol}_h=(2\pi)^{32}\prod_{k=1}^{N-1}\Ga(k/N)^{-c_k}=1125
  \qquad (N=15\ \text{и}\ 30),
\]
показатели $c_k$ точны (для $N=15$: $c_k=4$ при десяти $k$,
$c_k=6$ при $k=3,6,9,12$); остаток $1.01\cdot10^{-119}$;
$\det V^2=1265625/256$; дефляция PSLQ ($N=15$) восстанавливает
целочисленный вектор показателей; $N=30$ сертифицирован
синус-подсчётом.
\end{theorem}

\begin{proof}
Произведение Коула--Селберга по характерам: показатели собирают
кратности проводников $h_d$. Вычисление при $dps=120$ даёт $1125$ с
остатком $10^{-119}$; согласование с $\sqrt{\disc}$ из H1 точное.
Синус-подсчёт: $\Ga(k/30)\Ga((30-k)/30)=\pi/\sin(\pi k/30)$
замыкает тождество для чётных $k$ без PSLQ.
\end{proof}

\begin{theorem}[H3--H6]\label{th:H36}
\emph{(H3)} Каждый период Ферма --- $\Ga$-мономиал
$\frac1N\zeta^{ra+sb}\Ga(a/N)\Ga(b/N)/\Ga((a+b)/N)$ ($91/406$
периодов; квадратуры $10^{-121}$); уровень $30$ с нечётными $k$ не
сводим к $15$.
\emph{(H4)} Модули $|\omega(\tau_i)^2|$ --- $\Ga$-мономиалы;
фаза $\psi=\frac{(1+3\sqrt5)+i(\sqrt{15}-\sqrt3)}8=R/|R|$
восстановлена LLL; $|R|=3-\sqrt5$ --- единица $\QQ(\sqrt5)$.
\emph{(H5)} Лестница $\pi/7$, $\pi/15$, $\pi/30$ --- ступени одной
лестницы $\mu_N$-квантов; веса $4\sin(\pi k/N)$ согласованы.
\emph{(H6)} Граница знаменателей $Q_{\text{стенд}}=2N\max s_i=480$;
рационализация F замыкается без скана.
\end{theorem}

\begin{proof}
(H3) --- замкнутая форма периода + сертификат B ($69$ тестов,
$3.9\cdot10^{-36}$) + глубокие квадратуры ($10^{-121}$);
несводимость: $\Ga(1/30)$ не редуцируется инвентарём соотношений к
знаменателю $15$. (H4) --- LLL в $\QQ(\sqrt{-3},\sqrt5)$
восстанавливает точные координаты; $R\bar R=(3-\sqrt5)^2$ точное.
(H5) --- фазы $\pi/7$ (Клейн), $\pi/15$, $\pi/30$ --- уровни
конструкции $\delta=\pi/N$; веса согласованы с таблицами
Рамануджа. (H6) --- $2\cdot30\cdot8=480$; все измерения внутри
интервалов единственности F1.
\end{proof}

\begin{databox}[Клейн --- стенд в составе H]
Гладкость: $28(xyz)^3=0$; род $3$; $\Delta=-343=-7^3$; $c_4=105$;
$j=-3375=-15^3$; форма $W^2=v^3-35v-98$; корни $7,
\frac{-7\pm\sqrt{-7}}2$; $\Omega=1.9333117056\ldots$ (три схемы,
согласие $1.4\cdot10^{-32}$); $\tau=(1+\sqrt{-7})/2$;
$j(\tau)=-3375$ из периода ($57$ знаков, якорь $j(i)=1728$);
$\zeta+\zeta^2+\zeta^4=\tau-1$; изотипика $(1,1,0,1)$;
$\lambda_1$(Хивуд)$=\sqrt2$, кратность $6$; код Фано $C_7+\mu_4$.
\end{databox}

\begin{statusbox}[итог]
Сертификат H принят ($6/6$): циклотомическая решётка, объём
$1125$ в полной $\Ga$-нормировке, $\Ga$-мономиальные периоды,
Гросс-нормировка, лестница, знаменатели --- арка программы
замкнута. Время счёта ~2 минуты.
\end{statusbox}

\begin{verifybox}[как воспроизвести]
\texttt{python3 laboratory.py} $\to$ <<Сертификаты $\to$ H>>;
<<Стенды $\to$ Клейн>>; Lean: \texttt{stand\_disc},
\texttt{snf\_product}, \texttt{klein\_integers}.
\end{verifybox}

\vfill
{\color{linkblue}\hrule height 0.6pt}
\vspace{0.5em}
{\small\color{gray}
Лицензия: индивидуальная исключительная лицензия Исаева Исхака Хамзатовича.
Все права на вычисления, формулы и тексты принадлежат автору.
Верификация: \texttt{python3 laboratory.py} (пункт <<Стенды>>/<<Сертификаты>>),
Lean: \texttt{verification/lean/HodgeLaboratory.lean}.
}
\end{document}
